\subsection{Measurable functions defined on a measurable subset}

PROPOSITION 15. --- Let X be a locally compact space, \(\mu\) a measure on X, A a \(\mu\) -measurable subset of X, and f a mapping of A into a topological space F. The following conditions are equivalent:

\begin{enumerate}
\def\labelenumi{\alph{enumi})}
\item
  The set \(\mathfrak{H}\) of compact subsets \(K\) of \(A\) such that the restriction of \(f\) to \(K\) is continuous, is \(\mu\) -dense in \(A\) .
\item
  There exists a set \(\mathfrak{K}\) of compact subsets of A, \(\mu\) -dense in A, such that the restriction of \(f\) to every \(\mathrm{K} \in \mathfrak{K}\) is \(\mu_{\mathrm{K}}\) -measurable.
\item
  There exist a homeomorphism \(j\) of \(\mathrm{F}\) onto a subspace of a topological space \(\mathrm{G}\) and a \(\mu\) -measurable mapping \(g\) of \(\mathrm{X}\) into \(\mathrm{G}\) , such that \(g|_{\mathrm{A}} = j \circ f\) .
\item
  Every extension of f to a mapping of X into F, constant on X - A, is \(\mu\) -measurable.
\end{enumerate}

It is clear that \(a)\) implies \(b)\) and that \(d)\) implies \(c)\) . The fact that \(c)\) implies \(a)\) follows from condition \(c)\) of Prop. 12 of No.~8. On the other hand, \(b)\) implies \(a)\) : for, Def. 1 shows that, for each \(K \in \mathfrak{K}\) , the set of subsets \(H \in \mathfrak{H}\) contained in \(K\) is \(\mu_{K}\) -dense in \(K\) (No.~8, Prop. 12, \(c\) ), and Prop. 13 of No.~8 shows that \(\mathfrak{H}\) is \(\mu\) -dense in \(A\) . It remains to see that \(a)\) implies \(d)\) . Let \(g\) be an extension of \(f\) to \(X\) , constant on \(X - A\) . For every compact subset \(L\) of \(X\) , \(L \cap A\) and \(L \cap (X - A)\) are \(\mu\) -integrable; therefore, for every \(\varepsilon > 0\) , there exist a compact subset \(P \subset L \cap A\) and a compact subset \(Q \subset L \cap (X - A)\) such that

\[
| \mu | ((L \cap A) - P) \leqslant \varepsilon / 4 \quad \text { and } \quad | \mu | ((L \cap (X - A)) - Q) \leqslant \varepsilon / 4.
\]

On the other hand, there exists a set \(H \in \mathfrak{H}\) contained in P such that \(|\mu|(P - H) \leqslant \varepsilon/2\) ; the restriction of g to the compact set \(K = H \cup Q\) is then continuous (g being constant on Q) and \(|\mu|(L - K) \leqslant \varepsilon\) , which completes the proof.

DEFINITION 8. --- Let X be a locally compact space, \(\mu\) a measure on X, and A a \(\mu\) -measurable subset of X. A mapping f of A into a topological space F is said to be \(\mu\) -measurable if it satisfies the equivalent conditions of Prop. 15.

If A is locally \(\mu\) -negligible, then every mapping of A into F is therefore \(\mu\) -measurable.

COROLLARY 1. --- Let X be a locally compact space, \(\mu\) a measure on X, A a \(\mu\) -measurable subset of X, and \(f\) a \(\mu\) -measurable mapping of A into a topological space F. Let \(\mathfrak{K}\) be a set of compact subsets of X, \(\mu\) -dense in X. Then, there exists a partition of A formed by a locally negligible set N and a locally countable family \((\mathrm{K}_{\lambda})_{\lambda \in \mathrm{L}}\) of sets \(\mathrm{K}_{\lambda} \in \mathfrak{K}\) , such that \(f|_{\mathrm{K}_{\lambda}}\) is continuous for every \(\lambda \in \mathrm{L}\) .

In view of No.~9, Prop. 14, it suffices to show that the set \(\mathfrak{H} \subset \mathfrak{K}\) of subsets \(K \in \mathfrak{K}\) such that \(K \subset A\) and \(f|K\) is continuous, is \(\mu\) -dense in \(A\) . Now, it follows at once from Prop. 1 of No.~1 and condition \(d)\) of Prop. 15 that, for every compact subset \(K_0\) of \(A\) and every \(\varepsilon > 0\) , there exists a subset \(K \subset K_0\) belonging to \(\mathfrak{K}\) such that \(|\mu|(K_0 - K) \leqslant \varepsilon\) and \(f|K\) is continuous; the conclusion therefore follows from Prop. 12 of No.~8.

COROLLARY 2. --- Let K be a compact subspace of X; in order that a mapping f of K into a topological space F be \(\mu\) -measurable, it is necessary and sufficient that it be \(\mu_{K}\) -measurable.

In view of Lemma 2 of No.~7, this follows at once from Prop. 1 of No.~1 and condition \(a\) ) of Prop. 15.

PROPOSITION 16. --- Let A be a locally countable set of μ-measurable subsets of X and let \(B = \bigcup_{A \in \mathfrak{A}} A\) . For a mapping f of B into a topological space F to be μ-measurable, it is necessary and sufficient that the restriction of f to every \(A \in \mathfrak{A}\) be μ-measurable.

We have already observed (No.~9) that B is \(\mu\) -measurable. The condition being obviously necessary, let us prove that it is sufficient. Thus, let K be a compact subset of B. By hypothesis, there exists a sequence \((\mathbf{A}_{n})\) of sets belonging to \(\mathfrak{A}\) such that the \(K \cap A_{n}\) form a covering of K. Set \(C_{0} = K \cap A_{0}\) and

\[
\mathrm{C} _ {n} = \mathrm{K} \cap \mathrm{A} _ {n} \cap \mathbf {C} \left(\bigcup_ {i <   n} \mathrm{C} _ {i}\right)
\]

for n \textgreater{} 0, so that the nonempty \(C_{n}\) form a partition of K into \(\mu\) -integrable sets. Since the restriction of f to \(C_{n}\) is \(\mu\) -measurable, there exists a partition of \(C_{n}\) formed by a \(\mu\) -negligible set \(N_{n}\) and a sequence \((\mathrm{L}_{mn})_{m \geqslant 0}\) of compact sets such that \(f|L_{mn}\) is continuous. Since \(N = \bigcup_{n} N_{n}\) is \(\mu\) -negligible, we see that condition a) of Prop. 15 is satisfied, whence the proposition.

Property d) of Prop. 15 makes it possible to immediately generalize the properties of measurable functions defined on all of X, observed in Nos. 2 to 5, to measurable functions defined on a measurable subset A of X; these generalizations are left to the reader. We only point out explicitly that the principle of localization (No.~2, Prop. 4) remains valid when it is assumed that each of the functions \(g_{x}\) is only defined in \(V_{x}\) (or almost everywhere in \(V_{x}\) ) and is measurable.

